\section{Results}

\subsection{Metric definition}

Artifact 001 declared six finite compression metrics: support size, support
radius, mean distance to receipt home, maximum distance to receipt home,
support entropy, and concentration ratio. These metrics were defined as
descriptive finite-support quantities and were explicitly separated from any
physical gravity claim.

\subsection{Source discovery and priority}

The inherited project scan found 50 candidate compression sources across
Projects 19, 29, 30, and 31. The highest-priority inherited source was the
Project 19 holonomy return compression audit.

The G900 Observatory appliance scan found 444 candidate appliance sources.
Prioritization placed the return-cell one-step information transport support
audit and the six-nine return-home receipt ledger at the top of the appliance
source list.

\subsection{Return-cell metric result}

The admitted return-cell one-step information transport produced a positive
finite compression metric result under the conservative target-home proxy.

The observed shift was:

\begin{itemize}
  \item mean distance to receipt home: $2.75 \rightarrow 0.0$,
  \item support radius: $5 \rightarrow 0$,
  \item support size: $3 \rightarrow 3$,
  \item support entropy: $1.5 \rightarrow 1.5$.
\end{itemize}

Thus the return-cell transport did not shrink support. It aligned support
homeward.

\subsection{Home contrast}

The target-home proxy was compared against alternate finite home sets. The
target-home proxy produced exact after-arrival. The source-home set and the
six-nine pair did not preserve the compression signature. The full return-cell
slot set was exact but non-discriminating because it contained both source and
target support.

\subsection{Theorem verification}

The return-cell home-alignment theorem candidate was verified by dependency
check. The verifier checked 24 obligations and recorded zero failures.

The verified local statement is:

\begin{quote}
For the admitted return-cell one-step information transport, the source
support maps exactly into the admitted target-home slot set under the tested
finite slot-distance metric.
\end{quote}

\subsection{Replay stability}

Replay stability was tested under four equivalent presentations of the same
transport history: original order, reverse row order, sorted pairs, and unique
support. All four replay presentations preserved the compression signature.

\subsection{Perturbation profile}

The perturbation audit tested 24 cases. Of these, 23 preserved the compression
signature and 14 preserved exact-home arrival. The single signature failure
was the target rotation by five slots.

This shows that the return-cell compression signature is robust but not
absolute. It has a finite tolerance band.

\subsection{Return-cell checkpoint}

Artifact 009 recorded the first local Project 32 checkpoint:

\begin{quote}
A local replay-stable bounded homeward-alignment compression signature is
verified for the admitted return-cell one-step information transport.
\end{quote}

\subsection{Six-nine return-home result}

The six-nine return-home ledger passed as a strong receipt surface. It recorded
home arrival with receipts and ten route holds. However, the ledger was not
metric-ready in the same sense as the return-cell transport because it did not
expose a directed source-to-target slot history.

Subsequent bridge inspection found that the six-nine rowsets expose a surface
slotset anatomy. The dominant parsed slotset class was expanded surface with
gap, with slotset $\{3,6,9,11,12\}$ appearing 122 times.

\subsection{Surface taxonomy}

Artifact 016 recorded two distinct compression-relevant surface types:

\begin{itemize}
  \item return-cell transport: a metric-ready directed transport surface;
  \item six-nine return-home: a receipt-strong surface rowset candidate.
\end{itemize}

Artifact 017 checkpointed the phase with one verified local compression result
and one distinct six-nine surface compression candidate taxonomy.
